\begin{table}[htbp]
\centering
\caption{\new{Heterogeneity by RP-9 exposure in Legislature 56, dual outcome.}}
\label{tab:t2_dual}
\begin{threeparttable}
\small
\begin{tabular}{lcc}
\toprule
Subgroup & Gov. outcome & Centrão outcome \\
\midrule
RP-9 supporters ($d_{rp9}=1$) & $-0.48$ & $-0.07$ \\
Non-supporters ($d_{rp9}=0$) & $-2.27^{***}$ & $-0.93^{**}$ \\
\bottomrule
\end{tabular}
\begin{tablenotes}
\small
\item Notes: IV-DML estimates of the pork-for-votes effect (pp per R\$1M) in Legislature 56, split by RP-9 supporter status. A deputy is classified as RP-9 supporter if she appears as \emph{parlamentar solicitante} of at least one rapporteur transfer in 2020--2022 in the disclosure files released under ADPF 854 (3.2\% of deputy-vote observations). The estimates show that the visible-amendment backfire concentrates in the non-supporters; the effect is statistically indistinguishable from zero among supporters, consistent with substitution between visible (RP-6) and opaque (RP-9) channels.
\end{tablenotes}
\end{threeparttable}
\end{table}